% =====================================================================
% discussion.tex
% Content: why criterion adjustment typically dominates; the benefit/
% cost asymmetry and its biological reading; experimental-design
% guidance; new falsifiable predictions (anti-cue inversion, the
% decorrelation channel, conservation-form sensitivity); scoped
% limitations.
% =====================================================================

\section{Discussion}
\label{sec:discussion}

We set out to ask when value-directed attention matters: under what
conditions does re-allocating perceptual resources toward a high-value
location gain meaningfully over an observer that has already optimised
its decision criteria and its validity-driven attention? The model
returns a layered answer. Across a broad parameter space, criterion
adjustment is typically the dominant lever for encoding value, capturing
a median of roughly three-quarters of the reward gain
(Section~\ref{sec:results-criterion}). The advantage of value-directed
attention, where it exists, is non-monotonic in the benefit/cost ratio
$\Rsens$ and is governed by a closed-form lower edge $\rdagger(\val)$
(Section~\ref{sec:results-vda-nonmonotonic}); it is concentrated in a
graded corner of the design space rather than distributed across it
(Section~\ref{sec:results-graded}); and the direction of optimal
allocation never inverts when the cue is predictive, a conditional
theorem holding for $\valid \ge 1/\Nloc$
(Section~\ref{sec:results-noninversion}). Together these results
describe a task class in which the neural machinery for value-driven
re-allocation is, over much of the natural operating range, redundant
with a simpler criterion-based strategy --- and they delimit precisely
where it is not.

\subsection*{Why criterion adjustment is typically dominant}

Criterion adjustment and value-directed attention exploit the same cued
information by different routes, and the routes differ in price. Shifting
the decision threshold at a high-value location requires no perceptual
reallocation: a more liberal criterion there trades false alarms for
hits while leaving sensitivity at every other location untouched. Spatial
re-allocation, by contrast, pays a genuine perceptual cost --- gain
acquired at the cued location is taken from the others. When the cue is
predictive, validity-driven attention has already concentrated processing
where it is most likely to be needed, and the residual value information
is most economically booked into the decision thresholds rather than into
further perceptual gain. This is why, across the sweep, the criterion
fraction $\CF$ concentrates well above one-half (median $0.7552$ for the
value-coupled reward variant and $0.7682$ for the equal-reward
variant, both at the additive conservation rule), with most of the
distribution lying in $[0.30, 1.00]$
(Section~\ref{sec:results-criterion}). The lever does not dominate
everywhere: in the benefit-dominant, low-validity corner --- where
attentional enhancement is efficient and the cue is weakly informative
--- the criterion fraction falls and value-directed attention contributes
materially. The summary is therefore a central tendency with a
substantial tail across the parameter space.

\subsection*{The benefit/cost asymmetry and its biological reading}

The benefit/cost ratio $\Rsens$ is the model's compact summary of a
dissociation observed in cortex: attentional enhancement at attended
locations and suppression at unattended locations are separately
modulable, and the normalisation framework that captures single-unit
gain effects treats them as distinct operations rather than two ends of
one knob
\cite{reynolds_heeger2009_normalization, mcadams_maunsell1999_v4_tuning,
treue_martinez_trujillo1999_feature_attention, carrasco2011_visual_attention_25y}.
A large $\Rsens$ describes an enhancement-dominant regime in which
attending a location buys more sensitivity than is lost elsewhere; a
small $\Rsens$ describes a suppression- or cost-dominant regime in which
re-allocation is expensive. The non-monotonicity of the value-directed
advantage in $\Rsens$ (Section~\ref{sec:results-vda-nonmonotonic}) then
has a clean reading. When enhancement is cheap (large $\Rsens$), a
value-blind observer already attends the cued location and value adds
little; when enhancement is very expensive (small $\Rsens$), no observer
should re-allocate and again value adds little. The advantage peaks in
between, in a cost-dominant window where value-blind attention is too
expensive to deploy but the value differential makes targeted
re-allocation worthwhile --- exactly the band whose lower edge the
closed-form threshold $\rdagger(\val)$ marks
\eqref{eq:r-dagger}. The model thus locates value-directed attention not
where attention is easiest but where value is the deciding factor, which
is a more demanding and more diagnostic regime than the one many designs
target.

\subsection*{Guidance for experimental design}

These results carry directly into how one should design an experiment to
observe value-directed attention. Because the advantage is concentrated
in a graded corner --- low cue validity, high value contrast,
cost-dominant asymmetry, and low baseline sensitivity
(Section~\ref{sec:results-graded}) --- a design placed outside that
corner will tend to produce null value-directed effects even when the
underlying mechanism is intact, simply because optimal criterion
adjustment already captures nearly all of the available reward. Many
standard spatial-cueing paradigms, which use highly valid cues and
moderate value manipulations, operate in this dormant regime. The model
makes the boundary quantitative: at cue validity $\valid \ge 0.95$ the
predicted advantage sits at the grid floor for every asymmetry and
correlation we examined; the dormancy survives down to $\valid \ge 0.80$
when cross-location correlation is negligible but admits a small
correlation-conditional signal otherwise; and a $\valid \ge 0.75$
criterion is too permissive, because cost-dominant designs in
$\valid \in [0.60, 0.80)$ admit a peak advantage near $0.16$
(Section~\ref{sec:results-graded}). An experiment aiming to isolate
value-driven attention from criterion-based strategy should therefore
move deliberately into the cost-dominant, low-to-moderate-validity,
high-contrast corner rather than maximising cue validity.

\subsection*{New predictions}

The model makes three predictions that go beyond restating where the
advantage lives. First, although optimal allocation never inverts under a
predictive cue, it \emph{can} invert under a counter-predictive
``anti-cue'': when validity falls below chance, $\valid < 1/\Nloc$,
re-allocating attention \emph{away} from the nominally cued location
becomes optimal in a substantial fraction of cells
(Section~\ref{sec:results-noninversion}). The inversion is sharply
bounded in value contrast --- it is confined below the closed-form
boundary $\valid < 1/[(\Nloc-1)\val + 1]$ --- and is absent throughout
the predictive regime, which makes it a clean, falsifiable signature: an
observer who is genuinely re-allocating perceptual resources, rather than
only shifting criteria, should show inverted allocation under an
anti-predictive cue and never under a predictive one.

Second, the cross-location correlation $\corr$ is an active lever, not a
nuisance term, and its effect on the value-directed advantage changes
sign with the benefit/cost ratio. In cost-dominant designs ($\Rsens$
small) raising $\corr$ suppresses the advantage, whereas in
balanced-to-enhancement-dominant designs ($\Rsens \gtrsim 1$) it
amplifies it (Section~\ref{sec:results-graded}). The amplification can be
dramatic at design points that look dormant under independent noise: at
one cost-dominant high-contrast cell the advantage rises from a
negligible $0.0007$ at $\corr = 0$ to $0.0676$ at $\corr = 0.2$, roughly
a hundredfold increase driven entirely by admitting the empirically
typical level of inter-neuronal correlation
\cite{CohenMaunsell2009, RuffCohen2016, Srinath2021}. The prediction is
that whether a given design reveals value-directed attention can depend
on the correlation structure of the neural code, so that a paradigm that
is null under one correlation regime may not be under another.

Third, the position of the criterion/attention split is sensitive to the
form of the benefit/cost conservation rule. Replacing the additive rule
($\benefit + \cost = 2$) with the multiplicative endpoint
($\benefit \cdot \cost = 1$) leaves the central tendency of $\CF$ almost
unchanged --- the median moves by less than $0.005$ --- but it deepens
the tail: the fraction of cells with $\CF < 0.5$ roughly doubles, and the
per-cell criterion fraction never rises under the multiplicative rule
(Section~\ref{sec:results-criterion}). The dominance of criterion
adjustment is thus robust as a central claim but rule-dependent in its
tail, and the size of the attention-favourable tail is itself an
observable consequence of how enhancement and suppression trade off ---
a structural prediction about the resource constraint, not a free
parameter. The equal-reward convention sharpens the same point: its
deeper $\CF$ tail and its lower median inversion threshold
(Table~\ref{tab:noninv-tally}) indicate a somewhat larger
attention-favourable region and a somewhat higher anti-cue inversion
incidence than the value-coupled reward variant, so the reward
structure of a task is itself a design variable that shifts the regime
boundaries.

\subsection*{Scope and limitations}

Several boundaries delimit these statements. The benefit/cost
conservation constraint relating the benefit and cost weights
(Section~\ref{sec:model}, Equation~\eqref{eq:beta-gamma}) is one member
of a one-parameter family spanning the additive
($\benefit + \cost = 2$) and multiplicative ($\benefit \cdot \cost = 1$)
endpoints; the headline numbers are reported under the additive rule
with a band over the family, and which member corresponds to a
particular biological resource constraint is not fixed by the model. The single asymmetry ratio $\Rsens$ summarises a
homogeneous population; lifting it to a per-location vector and the uncued
allocation to the full $\Nloc$-dimensional simplex leaves the optima
essentially intact under modest heterogeneity, but the closed-form
dependence of the optimum on the spread of $\Rsens$ across locations is
not derived here. The decision rule treats per-location decision noise as
fixed; an attention-coupled noise channel $\sigma_i(\alpha)$ would, if it
were operative, enter alongside decorrelation as a further axis of the
decomposition and would shift the interpretation of the criterion
fraction in the cost-dominant regime, and we leave it as a natural
extension rather than including it. The transfer function relating
attention to sensitivity is taken from a family of monotonic forms; the
mechanism statements (Definition~\ref{def:three-levers},
Equation~\eqref{eq:r-dagger}) are stated for a generic transfer function
and do not depend on the particular family, though the empirical bands
do. The correlation lever assumes a single equicorrelated structure;
distance- or feature-dependent covariances break the one-factor
reduction and are out of scope. Finally, the model is normative and
stationary: it describes the optimal policy for an observer at a fixed
parameter setting, not how that policy is learned, and it makes no claim
about the neural implementation of any of the three levers beyond the
correspondence between the asymmetry ratio and the enhancement/suppression
operations of cortical gain control. Its contribution is a quantitative
map of when value-directed attention is the deciding mechanism and when an
optimally-set criterion suffices --- a map that a downstream empirical or
computational program can use to place its experiments where the
mechanism it cares about is actually at stake.
